\chapter{Why Did Socrates Keep Asking Questions?}
\section{A Bronze Plaque Bearing a Name}
Begin with something small: a thin bronze plaque scarcely larger than a bus card, its edges polished by wear, engraved with a person's name and a mark identifying his local community.

Many such plaques have been excavated in the Athenian agora. Archaeologists call them juror's tickets, or \textit{pinakia}. Over twenty-four centuries ago, Athens had no professional judges. Trials relied on lots: adult male citizens volunteered for jury service, and those selected received a plaque admitting them to court, where they heard the case and voted guilty or not guilty. Small as it was, it embodied an ordinary citizen's most tangible power. Judgement came not from king or priest but from hundreds of ordinary people---cobblers, carpenters, olive farmers.

One day in 399 BCE, about five hundred such plaques entered an Athenian court. The defendant was a seventy-year-old eccentric, famously ugly, barefoot year-round, wearing one robe through every season. He charged no fees, wrote no books, and spent his life stopping passers-by in the square to ask, ``What is justice?'' ``What is courage?'' ``Are you sure you know?'' Almost every Athenian knew him; comic theatre made him a joke.

That day, five hundred ordinary Athenians heard prosecution and defence, then voted. The majority found him guilty. He received a death sentence, and a few days later calmly drank hemlock in prison before his students and died.

This is our starting point, and the enormous question beneath the small bronze plaque: why did Athens, the city-state then proudest of its freedom of speech, vote to kill an old man who asked questions? How did questioning become a capital offence?

\section[Into the Fifth-Century Athenian Agora]{Entering the Athenian Agora in the Fifth Century BCE}
To understand, return to the scene.

The later fifth century BCE, in Athens's agora. Enter first through your nose: olive oil and grilled fish mingle with the sweat of thousands. Enslaved people squeeze past carrying water jars, fig sellers shout, stonecutters ring out at a distant building site. Athens has recently won a great victory and is sliding towards a great war, overflowing with wealth and burning with tension.

A crowd gathers beneath a colonnade. At its centre, an elegantly dressed foreigner in fine cloth speaks with an actor's modulation. He is a fee-charging teacher of the most sought-after skill: speaking. How to persuade thousands in the assembly to vote for you; how to turn black into grey and grey into white in court. His pupils are wealthy Athenians' sons, and fees are high. Athens has more than a dozen such teachers. People call them \textit{sophistai}, Sophists: skilled possessors of wisdom.

Nearby, another circle sits in a column's shade. At its centre is the ugly old man, Socrates. He takes no payment and has no script. He is questioning a young aristocrat fresh from a Sophist's lesson:

``You want to become a just person. What, then, is justice?''

``Justice is \,\ldots\, obeying the law,'' the young man replies.

``Good. Spartan law requires weak newborns to be abandoned in the mountains. Is obeying that law just?''

The young man frowns. ``Well \,\ldots''

``Or suppose your general orders you to retreat, abandoning a wounded comrade. Is obeying just, or disobeying?''

Onlookers laugh. The young man blushes. The old man looks sincere, even awkward, as though genuinely seeking instruction. After half an hour, however, everything the youth believed about justice lies in pieces, and he cannot produce one complete answer.

This was an everyday Athenian spectacle: Sophists charged to teach victory; an old man freely forced admissions of ignorance. Notice that both questioned. One route yielded wealth and reputation, the other a death sentence. Where was the difference? That is what we shall examine.

\section{What Was at Stake in the Dispute?}
Before choosing sides, see the problem clearly.

On the surface, money and presentation: Sophists charged, Socrates did not; they lived comfortably, he went barefoot. Beneath lay three real questions.

First: \textbf{can virtue be taught?} Sophists said yes. Justice, courage, and civic government were teachable skills like navigation or medicine: pay, attend, graduate, apply. It sounds obvious; today's education industry rests on this conviction. Yet Socrates fixated on something odd. Athens's greatest and supposedly bravest, most just statesmen often had mediocre or wasteful sons. If virtue was a skill like medicine, why could a great physician's son not learn it? The question was shrewd: either virtue was no such skill, or Sophists took payment without teaching anything real.

Second: \textbf{where is the measure of right and wrong?} Plato's \textit{Theaetetus} records Protagoras, the most famous Sophist, saying, ``Man is the measure of all things.'' Broadly, wind is cold to whoever feels cold and warm to whoever feels warm. No single truth hides behind appearances, only particular people's experiences and judgements. Common sense today, it was explosive then. If each person and city has a measure, is justice whatever Athens declares? Does it change across borders? Socrates opposed this instability throughout his life. Behind courage, justice, and piety, he insisted, there should be something testable that does not sway with the crowd. Yet a lifetime's search produced no definition. Embarrassingly, a man convinced of standard answers died without supplying one.

Third, the question that brought him to court: \textbf{can a city-state tolerate endless questioning?} Socrates believed it a service. God had sent him as a gadfly to sting Athens, a large, lazy thoroughbred, awake. Athenians heard something else: this man taught youths to contradict fathers, mock tradition, doubt civic gods and democracy itself. The gadfly claims to heal the horse; the horse only feels pain. Who is right?

Together, these are three faces of one question: when does a society's ``why'' build, and when does it demolish? May questioners pursue it to the end? May society demand that they stop?

Remember your present inclination: old man or horse? Do not settle yet. We must let every side make its strongest case.

\section{Give Each Side Its Full Case}
Histories of Socrates easily choose automatically: a martyr for truth against a mob murdering wisdom. Resist that conclusion for now. This chapter practises making each position so strong that its supporters cannot object. That is steelmanning, opposite to straw-manning, which weakens and ridicules a position before knocking it down. First build the strongest version; then decide whether to agree.

\textbf{First, the Sophists.}

Textbooks often cast them as villains selling sophistry and confusing right with wrong. Their case is strong. First, they broadened educational access. Before them, civic leadership skills passed within aristocratic families. They placed success in public life on sale to anyone who could pay. Newly wealthy commoners could buy what ancestry had monopolised. Second, they served democracy. Public offices were allotted, hundreds of citizens voted in court, and assembly speeches determined policy. A citizen unable to speak was half a citizen. Teaching speech fuelled democracy. Third, and deepest, ``Man is the measure'' might express humility rather than cynicism. Where oracles, tradition, and noble authority monopolised correct answers, locating judgement in particular people brought the right to think down from the altar. The Enlightenment would do the same two thousand years later.

\textbf{Next, the Athenians and the accusers.}

This is harder and more important. Defending the opponents of a dead hero wins little popularity. Listen seriously.

A middle-aged ordinary citizen who has lived through war speaks. You see us killing an inquisitive old man, he says, but not what we have just escaped. In 404 BCE, Athens lost the great war against Sparta. Thirty aristocrats took power, the Thirty Tyrants. In eight months they executed about fifteen hundred citizens without trial, confiscated property, and expelled opponents. In a city-state of tens of thousands, this terror touched everyone. Their leader Critias had studied in Socrates's circle. So had Alcibiades, the beautiful genius and notorious traitor who defected first to Sparta, then Persia, nearly destroying his homeland.

``I am not saying he taught them murder,'' the citizen continues. ``But tell me: someone spends his life teaching youths to dismantle reverence for gods, laws, and fathers. His two outstanding associates become a murderous tyrant and a traitorous general. We have crawled from the bodies and painfully rebuilt democracy. Should we not be afraid?''

A technical detail matters. Restored democracy proclaimed an amnesty forbidding prosecution for earlier political offences. Nobody could prosecute Socrates because his student was a tyrant. Hence the 399 BCE indictment listed other charges: rejecting the city's gods and introducing new divinities, and corrupting youth. Call these pretexts if you wish. But where religion and civic life intertwined, with sacrifices before meetings and campaigns, disbelief in civic gods was not a private preference. It resembled rejecting the community's fundamental rules.

The citizen's case has holes: fear is not evidence, association not guilt. Yet first acknowledge that this is not a band of savages burning witches. It is a gravely wounded democracy using legal procedures to handle a real fear clumsily. Only then do the following questions acquire weight.

\textbf{Finally, voices beyond Socrates's own circle.}

Questioning was not Athens's patent. Around the same era, Indian thinkers challenged one another in Ganges forests and towns, later developing rigorous debating rules: thesis, reason, example, inference, one wrong step losing everything. This learning was called \textit{hetuvidya}, the science of reasons. Centuries later, Xuanzang studying at Nalanda still recorded a culture judging learning through debate. Jewish rabbis argued over scripture for more than a millennium, their records forming the \textit{Talmud}. Distinctively, it carefully preserves minority views that lost. A traditional saying holds, broadly, that in disputes seeking truth, even the defeated side deserves preservation. In China, ``Lesser Selection'' in the \textit{Mozi} opens:

\begin{quote}
Disputation clarifies the division between right and wrong, examines the principles of order and disorder, distinguishes sameness from difference, and investigates the relation of names to realities.
\end{quote}
In essence, debate distinguishes right and wrong, sameness and difference, and tests whether words match things. Its purpose is clarity, not victory.

Together, these suggest a cross-cultural discovery: almost every civilisation independently invented the questioning machine. They also discovered its failures: it hurts, unsettles order, and can be captured by competitiveness. Some installed safeguards, such as preserving losing opinions. Others, like Athens, executed the machine's inventor.

\section[Thinking Bridge: Test a Definition]{Thinking Bridge: Three Counterexamples against a Definition}
Detach Socrates's signature skill and put it in your pocket. It costs little and has three steps:

\textbf{First, request a definition.} Hearing a large word---courage, fairness, freedom, maturity, effort---neither agree nor disagree. Ask what it means and request a criterion.

\textbf{Second, find a counterexample.} Something that meets the definition but intuitively should not count, or fails it but intuitively should. One suffices; two are better.

\textbf{Third, revise and return to Step Two.} Continue until revision stops or you admit no definition is available.

Watch it work. Plato's \textit{Euthyphro} records a real dismantling, broadly as follows. At court Socrates meets Euthyphro, who is prosecuting his father and considers himself pious. Socrates immediately seeks instruction: what is piety? First answer: doing what I am doing, prosecuting wrongdoing even by a father. Socrates objects: an example, not a definition. Second: what gods love. Do gods quarrel? If so, one loves what another hates; is the same act pious or not? Euthyphro revises: what all gods love. Socrates asks whether they love it because it is pious, or it is pious because they love it. A tongue-twister with a sharp edge. The latter reduces piety to divine mood, without its own reason. The former means divine approval explains nothing about piety itself. At the end Euthyphro announces urgent business and leaves. Such is elenchus, refutation through questioning: without lifting a finger, make someone discover they cannot explain a word they use daily.

Other civilisations mirror it. In ``Autumn Floods'' in the \textit{Zhuangzi}, Zhuangzi and Hui Shi watch fish from a Hao River bridge. Zhuangzi says their unhurried swimming shows happiness. Hui Shi asks, ``You are not a fish; how do you know its happiness?'' Zhuangzi refuses a definition and turns the challenge back: ``You are not me; how do you know I do not know?'' Neither persuades the other, but their ancient exchange exposes a continuing psychological and philosophical difficulty: what warrants knowledge of another's feelings?

Three cautions govern this tool. First, winning is not its purpose. Socrates called his method midwifery. His mother delivered babies; he delivered others' ideas, testing them before birth and letting unsustainable ones miscarry. Use counterexamples merely to embarrass and the midwife becomes a provocateur. Second, counterexamples must be fair. Pointless contrariness says, ``Water quenches thirst? Try seawater.'' A testing counterexample strikes the definition's central weakness. Third, hardest of all, a method dismantling others' definitions dismantles yours too. Socrates was particularly resented because he exempted not even himself. Knowing one's ignorance must begin with oneself.

Finally, a term: irony, \textit{eironeia}, here means neither mockery nor sarcasm. Socrates's characteristic pose is professed ignorance: I know nothing; teach me. Then his requests for instruction dismantle the teacher's knowledge. This advance through apparent retreat became Socratic irony. It is gentle attack, or attacking gentleness. Few recipients found it gentle.

\section[In Court: Reading the Apology]{What He Said in Court: Reading the Apology}
Now read primary material. Its provenance comes first, because that is itself an evidence lesson.

Socrates wrote nothing. We know him mainly through two pupils: Xenophon, a soldier writing plainly, and Plato, a philosopher writing brilliantly. Brilliance begins the problem, because Plato later placed much of his own thought in his teacher's mouth. Historians call this the Socratic problem: we receive not the man but others' remembered and reworked versions. Thus cite ``Socrates as recorded in Plato's \textit{Apology},'' not simply ``Socrates said.'' This is no pedantry. It embodies four distinctions: what happened, how participants understood it, how recorders shaped it, and how we read it.

The \textit{Apology} records his defence. He tells, broadly, how his friend Chaerephon visited Delphi and asked whether anyone was wiser than Socrates. The priestess's oracle answered no. Socrates initially objected: I know I understand nothing; how can the god say this? Greeks believed Delphi spoke for Apollo, who could not lie. He therefore sought to refute the oracle by visiting reputedly wise politicians, poets, and craftsmen. Find one who truly understood and the oracle would fail.

The result disturbed him. Politicians spoke fluently of government but faltered under questions. Poets created enduring verses without explaining them, as though inspired rather than understanding. Craftsmen genuinely knew their crafts but mistook one skill for knowledge of other matters, an error overshadowing their real ability. His interpretation followed: the god did not mean Socrates knew much, but that the wisest person recognises their wisdom's worthlessness. Others did not know their ignorance; he at least did.

This is the source of ``knowing that you do not know.'' Notice its shape: not polite modesty, but a cold cognitive conclusion. Knowing ignorance's boundary is itself the highest human wisdom. At roughly the same time across the continent, Confucius told Zilu something strikingly parallel:

\begin{quote}
To know what you know, and recognise what you do not know: that is knowing. (\textit{Analects}, ``On Government'')
\end{quote}
Know when you know, acknowledge when you do not: this is genuine knowledge. Two teachers who never met, a continent apart, reached one conclusion: learning begins by honestly marking ignorance's boundary.

The court heard a still more famous declaration. Called a pest, Socrates accepted the image with one alteration, broadly: the city is a fine, large but sluggish horse, and I am the god-sent gadfly pursuing, provoking, and keeping it awake. Kill me easily, but you will struggle to find another. He added a sentence repeated countless times:

\begin{quote}
The unexamined life is not worth living. (Plato's \textit{Apology}, conventional Chinese rendering)
\end{quote}
Broadly: a life never inspecting itself does not deserve to be called human. Notice the setting. This is no salon maxim but a defendant addressing judges before a death sentence. Not encouragement to improve, but a testament: he knows it cannot save him and will only inflame matters, yet says it. Why? The next section compares answers.

\section[Why Socrates Had to Die]{Why Did Socrates Have to Die? Three Explanations}
The five hundred plaques have voted; the old man is dead. Evidence broadly establishes what happened, and we have heard prosecution and defence understandings. The hardest layer remains: why? Historians offer different autopsies of the same body.

\textbf{Explanation One: Political aftermath---a reckoning dressed in religious charges.}

The timing fits. The Thirty's terror ended in 403 BCE; trial came in 399, wounds still bleeding. Socrates taught their leader and befriended a traitorous general, while amnesty barred political prosecution. Impiety and corrupting youth supplied the available legal handles. A leading mover, Anytus, helped restore democracy. Here the death is a civil-war aftereffect: democracy kills not a philosopher but the unreachable tyrants' shadow.

This explains timing and charges, but leaves a blind spot. If the whole city wanted him dead, why was the vote so close? Later tallies differ, but describe a majority rather than a landslide. A slightly conciliatory defence and accepted penalty probably could have saved him under Athenian practice. A purely political reckoning should not yield so easily to a few soft words.

\textbf{Explanation Two: He really offended Athenian law---the jury was sincere, not manipulated.}

Set modern religion aside temporarily. Athenian gods were not private beliefs but enrolled members of the polity: festivals, oaths, campaigns, and trials operated in their names. Publicly claiming guidance from a mysterious divine voice while daily testing traditional beliefs sounded like demolishing load-bearing walls. Corrupting youth was not entirely invented either: he taught methods for dismantling paternal authority and civic common sense. The method was neutral, but he supplied no rebuilding plan for the ruins. Jurors therefore chose between real values, placing communal cohesion above one old man's right to question.

This restores Athenian dignity, at a cost. If we set aside the truths about the two tyrannical periods, whom exactly did he corrupt, and on what evidence? The indictment remains vague. Handling real fear through vague charges is a classic fracture in procedural justice.

\textbf{Explanation Three: He sought death---the trial was his final lesson.}

His defence seems unconventional everywhere. Athenian defendants wept, begged, and displayed children for sympathy. He did none of these, instead proposing punishment as free meals in the public hall, an Olympic victor's honour. After the death sentence, procedure allowed him to propose an alternative. A respectable fine was within friends' means; he initially refused, then reluctantly suggested an amount. In prison, Crito arranged an escape, but he rejected it: a lifelong teacher of lawfulness could not escape through lawbreaking (Plato's \textit{Crito}). Together, these suggest a reading: at seventy, why die ill in bed when death could become a final lesson, making Athens witness the quality of its laws, democracy, and justice by executing him?

This is the most moving literary explanation and the most dangerous. It romanticises a complex public event into one genius's solo performance, reducing five hundred jurors to props. It also fails to explain why he refuted each charge so earnestly. Someone simply seeking death need not argue so hard.

Each illuminates a corner and leaves shadow. That is honesty, not fence-sitting. Evidence for a twenty-four-century-old verdict is limited; anyone claiming one certain explanation has probably discarded some of it. Choose or combine, but state the evidence by which you exclude the others.

\section[Further Challenge: Defining Concepts]{Further Challenge: Do Concepts Have Standard Answers?}
Socrates sought definitions throughout life without obtaining one: courage, justice, piety, beauty, each demolished by counterexamples. Was his method flawed, or do these words \textbf{have no such definitions}?

Over two millennia later, Austrian philosopher Ludwig Wittgenstein offered an answer that kept philosophers awake. In his later \textit{Philosophical Investigations}, he asks readers to set theory aside and look at activities called games, \textit{Spiel}: board games, cards, ball games, children's play. What do they share? Chess lacks physical contest, football a board, solitary string figures an opponent, playing house even winning and losing. No feature belongs to every game and only games. Instead there is an overlapping network of likenesses, like a family: the eldest has Dad's eyes, the second Mum's nose, the third the eldest's chin, without one trait shared by all. He calls this \textbf{family resemblance}.

The idea's force is that it explains Socrates's repeated defeats. If courage, justice, games, and art gather examples through family resemblance rather than enclosing territory with a precise definition, defining justice against every possible counterexample was impossible from the outset. The thinker is not foolish; the assigned task does not exist. Counterexamples remain available because conceptual edges are fuzzy.

Do not declare victory yet. An easily misjudged counterexample: mathematical definitions withstand challenges. ``A triangle has three sides,'' ``A prime is divisible only by one and itself.'' Search the world and find neither a four-sided triangle nor a prime divisible by seven. What follows? Not that all definitions fail. Some are rock-solid; others always leak. The real question becomes why ethical and mathematical terms differ. One answer: mathematical concepts form constructed closed systems with fixed rules, whereas courage and justice grow from changing human practice and remain blurred. Another defends Socrates: precisely because boundaries blur, generations must keep questioning, not to reach an endpoint but to prevent unattended concepts from quietly changing shape. Both answers live on in linguistic and cognitive-scientific debate.

Take away a perspective, not a verdict. Next time anyone, including this book, strikes you with a categorical definition, ask whether the concept resembles a triangle or a game. If a game, its definer is either simplifying or staking a territorial claim.

\section{Today's Gadflies and Courts}
This ancient trial repeats around you weekly, with something else replacing poison.

First, the classroom. A student asks, ``You say the rule benefits us. Could you test it with a counterexample?'' Gadfly or provocateur? The tool marks the distinction by purpose. Midwife-like questioning seeks clarity and accepts scrutiny of itself. Contrarian questioning seeks another's defeat and its own victory, dismantling others but never itself. One sentence, two spirits. Athens's cautionary error was not suspecting questioners but failing to distinguish them, swatting gadfly and horsefly together.

Next, the internet: history's largest square and court. Someone speaks; hundreds of comment-section jurors vote; conviction takes half a day. Every Athenian role is online: paid coaches teaching verbal victory, the Sophists' descendants; irritating persistent questioners; frightened crowds convicting under vague accusations of suspect positions or impure motives. The procedure has not changed in twenty-four centuries, only accelerated ten-thousandfold. Athens at least allowed its defendant a whole day's speech.

The deepest echo is our opening tension: \textbf{how does challenging authority relate to public responsibility?} Socrates offered a difficult answer. He spent his life undermining Athens, mocking politicians, questioning beliefs, deriding democratic votes. Yet sentenced himself, he refused escape. Living there and enjoying legal protection, he argued, meant contracting with its laws. He could criticise the verdict and seek persuasion, but not tear up the contract upon losing (broadly, the \textit{Crito}). Together these halves make a complete Socrates: an unsparing critic and a citizen refusing to evade payment. Criticising a community and abandoning it differ fundamentally. The former is responsibility's highest form, because someone escaping need not make the community better.

Today the problem wears many clothes. Does identifying school problems smear the school? Does criticising home betray one's roots? Is challenging a company disloyal? Each generation must ask again: who needs whom more, horse or gadfly?

A new situation was unavailable to Socrates. Answers grow cheaper: machines supply any homework answer in three seconds. Good questions grow more valuable. Someone who detects an apology's assumptions or first asks what a trending question presupposes holds Socrates's old tools. The tools are unchanged; people able to use them seem, if anything, fewer.

\section[Your Turn: One of Five Hundred Plaques]{Your Turn: One of Five Hundred Bronze Plaques}
Return to the opening plaque.

Imagine being an Athenian juror in 399 BCE: cobbler, carpenter, or olive farmer. You endured war; a neighbour's property was seized under the Thirty. You heard agora debates, perhaps suffered the old man's questioning yourself. Today you heard the accusations and his unyielding, almost arrogant defence: god-sent gadfly, unexamined life unworthy, punishment as free public meals.

Now you must cast your vote.

Guilty or not guilty?

Before voting, weigh everything: the Sophists' case for education, democracy, and each person's right to think; the accusers' real fear, real bodies, communal self-preservation; the old man's honest ignorance, valuable questioning, and possibility of criticism joined with faithful citizenship; and every uncertainty left by the three explanations.

Go deeper, because Socrates would ask: what grounds your vote? A judgement of guilt or a feeling of annoyance? Are you sure you can explain the difference?

There is no standard answer. The five hundred Athenians had none either, only one day and two voting tokens. Your sole advantage is twenty-four centuries of hindsight. History's cruelest warning is that hindsight does not guarantee you would have voted more wisely in their place.
